\section{Datasets and Acoustic Characteristics}
\label{sec:data_alignment}\label{sec:acoustic_characteristics}

\noindent\textbf{Datasets and Annotations. }
ICBHI and SPRSound (SPR in tables) provide training data with different annotation units and label detail. ICBHI labels respiratory cycles as Normal (N), Crackle (C), Wheeze (W), or Both (C and W)~\cite{rocha2019icbhi}. SPRSound provides seven event categories, including Fine/Coarse Crackle and Wheeze, while its task studied here predicts Normal or Adventitious (abnormal)~\cite{zhang2022sprsound}. Both datasets provide patient identifiers and annotation boundaries.

HF labels include inhalation (I), exhalation (E), discontinuous adventitious sounds (D), Wheeze, Rhonchi, and Stridor~\cite{hsu2021hflung}. HF train also supplies auxiliary C/W supervision in LSAA w/HF; fixed-model evaluation on HF test distinguishes recordings with Wheeze, Rhonchi, or Stridor annotations from other eligible recordings.
KAUH provides 112 patients with 336 recordings (Bell, Diaphragm, and Extended filter versions grouped by patient)~\cite{fraiwan2021kauh} leaving 86 patients(258 recordings) for external evaluation after we mapped N, C/IC, EW/IEW, and ICEW to Normal, Crackle, Wheeze, and Both. Figure~\ref{fig:dataset_characteristics}(a) reports original label counts and data partitions.

\noindent\textbf{Acoustic Characteristics. }
%The shared C/W labels in ICBHI and SPRSound provide a basis for joint supervision. 
We use mono 16-kHz, 5-s inputs with repeat padding or front cropping, boundary fading, or consecutive windows for each sample and then produce spectral centroid, bandwidth, and peak frequency from Welch spectra~\cite{welch1967spectra}. Descriptors are aggregated by median within recording and then within patient (recording date for HF). We also conduct a log transformation and joint standardization for principal component analysis (PCA). 

Panel (a) shows that Normal is the largest category in both core datasets, despite their different label inventories. PCA shows partially overlapping dataset distributions (b); we therefore compare corresponding sound labels in (c) to examine whether the differences persist within sound categories. After grouping SPRSound's Fine/Coarse Crackle as Crackle and Wheeze+Crackle as Both, panel (c) shows lower median spectral centroids for ICBHI than SPRSound in all four corresponding groups. HF groups reflect D/W annotations, not confirmed sound absence, so no Normal group is assigned; KAUH codes N, C/IC, EW/IEW, and ICEW map to Normal, Crackle, Wheeze, and Both, with other codes excluded. These observations motivate testing whether joint attribute supervision supports both native tasks across the observed differences, using single-source comparisons and fixed-model evaluation on HF test and KAUH (Section~\ref{sec:evaluation}).

